\section{Conclusions}
\label{sec:conclusions}

The \textit{Kitsurai} project was developed with the goal of creating a distributed key-value store capable of offering predictable and guaranteed performance in multi-tenant environments, a requirement often overlooked in major existing open-source solutions.

The experimental results confirm the validity of the approach: the system manages to maintain controlled performance even under high load and in the presence of malicious clients, ensuring effective isolation between tenants and guaranteeing each a minimum share of bandwidth and an acceptable maximum latency.

\subsection{Limitations}
\label{subsec:limitations}

Despite the promising results, the current implementation has some limitations that provide opportunities for future work. In particular, the system still lacks:

\begin{itemize}
    \item An access control model for fine-grained permission management.
    \item A catch-up protocol for reintegrating nodes after prolonged outages.
    \item Robust mechanisms for conflict resolution during read operations.
    \item The ability for users to assign custom names to tables.
    \item Strategies for optimal capacity management of nodes, especially when tables have unbalanced requirements in terms of bandwidth or storage space.
    \item Support for dynamic addition and removal of nodes from the system.
    \item A more performant and encrypted RPC protocol based on existing technologies, for example QUIC~\cite{QUIC}.
\end{itemize}
